\scanpage{197}

% B01 / text / top, corollary 2 and NB
\mspage{101}\uline{Corollaire 2}\quad Toute structure \deleted{orientée} de contour $C$
sur $E$, \unsure{invarie|invariée} par $\Lambda_E$, admet exactement
deux orientations invariées par $\Lambda_E$, elles sont
opposées l'une à l'autre. De cette façon, l'ens
$\unsure{\mathcal C|\mathscr C}(E)$ de ces contours s'identifie à l'ensemble
quotient
% slanted
\margin{\uline{NB} Le sous groupe \deleted{\unsure{$D_E$ de}} $\unsure{\mathfrak S|\mathfrak G}_E$ qui normalise $\Lambda_E$ est le ss-groupe affine $\mathrm{Aff}_{\Lambda_E}(E)$, et $D_E^+$ se décrit comme l'image inverse de $\{\pm1\}\subset(\mathbb Z/n\mathbb Z)^*$ dans cette ext. de $(\mathbb Z/n\mathbb Z)^*$ par $\Lambda_E$. Les structures de contour envisagées invariées par $\Lambda_E$ sont alors $\overline{\mathrm m}$ invariées par $D_E$ (mais non par $\mathrm{Aff}(E)$ tout entier, bien sûr). Le groupe $D_E$ est alors simplement transitif sur l'ens des répères de toute structure \deleted{\unsure{polyg}} de contour envisagée (à l'exclusion des deux dégénérées)}

% B02 / display / upper middle, (22)
\[
 \mathcal C(E)\simeq\bigl(\Lambda_E\smallsetminus\{0,\overset{\inserted{\text{si }n\text{ pair}}}{\sigma}\}\bigr)/\mathord{\pm}\mathrm{id}_{\Lambda_E} \tag*{(22)}
\]

% B03 / text / middle, parenthesis
(\deleted{On} \deleted{prend} l'ensemble $\Lambda_E/\mathord{\pm}\mathrm{id}_{\Lambda_E}$ plus sympathique,
si on admet les deux contours invariés ``impropres'',
correspondants à $0$ et $\sigma$ $\ldots$) dans ce cas,

% B04 / text / middle, monoid and cancellations
\deleted{\unsure{$\Lambda(\mathbb Z)$}}\deleted{le groupe} \inserted{le monoïde} $\Lambda_n^\times=\mathbb Z/n\mathbb Z$ avec la multiplication
opère sur cet ensemble, qui \deleted{\unsure{doit admet une|admet un}}\inserted{est \deleted{\unsure{ainsi}}} sera
\deleted{pour} \unsure{non can.$^{\mathrm t}$|non can.|non com.$^{\mathrm t}$} \unsure{à|;} $\Lambda_n^\times$ comme ens à opérateurs \unsure{?}

% B05 / text / lower, polygonal structures
L'ens des structures \deleted{de} polygônales sur $E$
invariées par $E$ \deleted{forme un torseur}\inserted{est en corresp 1--1} avec
$\Lambda_E^*/\mathord{\pm}1$, et forme un torseur sous \deleted{\unsure{$\Lambda'_n=\Lambda$}}

% B06 / display / lower, (23)
\[
 \unsure{Z'_n|\mathbb Z'_n}=(\mathbb Z/n\mathbb Z)^*/\mathord{\pm}1\ . \tag*{(23)}
\]

% B07 / text / foot, operators
Pour $\unsure{\nu|v}\in(\mathbb Z/n\mathbb Z)^*$ resp $\unsure{\dot\nu|\dot v}\in Z'_n=(\mathbb Z/n\mathbb Z)^{\unsure{*|1}}/\mathord{\pm}1$,
je désigne par $T_\nu$ resp $T_{\dot\nu}$ l'opération
correspondante sur l'ens $\overrightarrow{\mathcal C}_E$ resp. $\widehat{\mathcal C}_E$ --
et cette \unsure{dénotation|denotation|opération} \unsure{les les identiant|les les identifiant|des les identiant|les éléments} pourra
être prolongée au cas de $\nu\in\mathbb Z/n\mathbb Z^{\unsure{*|\vphantom{1}}}$ (sans $*$ !)
$\dot\nu\in(\mathbb Z/n\mathbb Z)/\mathord{\pm}1$ --

% archive sheet number: 196
